\documentclass{beamer}
\usepackage{hyperref}
\usepackage[T1]{fontenc}
\usepackage{cite}
\usepackage{relsize}
\usepackage{animate}


\usefonttheme[onlymath]{serif}
%C++ symbol
\newcommand{\Rplus}{\protect\hspace{-.1em}\protect\raisebox{.35ex}{\smaller{\smaller\textbf{+}}}}
\newcommand{\Cpp}{\mbox{C\Rplus\Rplus}\hspace{3pt}}
%Scientific notation
\newcommand{\expnumber}[2]{{#1}\mathrm{e}{#2}}

\usepackage{latexsym,amsmath,xcolor,multicol,booktabs,calligra}
\usepackage{graphicx,tikz,listings,stackengine}
\usetikzlibrary{shapes,snakes,matrix, positioning,arrows.meta}
\setbeamertemplate{bibliography item}[text]

\author{Rafael De Luna Loredo}
\title{Programaci\'on en Python y \Cpp}
\subtitle{Programaci\'on Orientada a Objetos}
\institute{Facultad de Ciencias, UASLP}
\date{}
\usepackage{NJFU}

% defs
\def\cmd#1{\texttt{\color{red}\footnotesize $\backslash$#1}}
\def\env#1{\texttt{\color{blue}\footnotesize #1}}
\definecolor{deepblue}{rgb}{0,0,0.5}
\definecolor{deepred}{rgb}{0.6,0,0}
\definecolor{deepgreen}{rgb}{0,0.5,0}
\definecolor{halfgray}{gray}{0.55}

\begin{document}
	
	\AtBeginSubsection[]
	{
		\begin{frame}
			\begin{multicols}{3}
				\tableofcontents[currentsection,subsectionstyle=show/shaded/hide,subsubsectionstyle=show/shaded/hide]
			\end{multicols}
		\end{frame}
	}
	\AtBeginSubsubsection[]{
		\begin{frame}
			\begin{multicols}{3}
				\tableofcontents[currentsection,subsectionstyle=show/shaded/hide,subsubsectionstyle=show/shaded/hide]
			\end{multicols}
		\end{frame}
	}
	
	\AtBeginSection[]
	{
		\begin{frame}
			\begin{multicols*}{3}
				\tableofcontents[currentsection,subsectionstyle=show/shaded/hide,subsubsectionstyle=show/shaded/hide]
			\end{multicols*}
		\end{frame}
	}
	
	\begin{frame}
		\titlepage
		\begin{figure}[htpb]
			\begin{center}
				\includegraphics[width=0.2\linewidth]{fig/logo.png}
			\end{center}
		\end{figure}
	\end{frame}
	
	\begin{frame}{Indice}
		\begin{multicols*}{3}
			\tableofcontents[sectionstyle=show,subsectionstyle=show/shaded/hide,subsubsectionstyle=show/shaded/hide]
		\end{multicols*}
	\end{frame}


\section{Introducci\'on}

\begin{frame}{Programaci\'on Orientada a Objetos}
	Es un paradigma de programaci\'on que nos permite representar mejor conceptos de la realidad mediante objetos, que no es mas que una manera de llamar
	a una especie de estructura de datos pero que no trabaja como tal solo con datos sino con conceptos propios de un objeto real como un carro, como puede ser 
	el estado de la bater\'ia, nivel del tanque, la velocidad actual o velocidad m\'axima, de igual manera no solo estados sino tambi\'en acciones como puede ser
	el avanzar, pitar, encender o abrir cajuela. 
	Entonces la POO es una manera de interpretar objetos del mundo real en t\'erminos de programaci\'on.
\end{frame}

\begin{frame}{Elementos clave}
	\begin{itemize}
		\item Las acciones o comportamientos se llaman m\'etodos
		\item Se basan en una plantilla llamada clase
		\item Los datos o informaci\'on que contiene se llaman atributos 
		\item  Los m\'etodos de una clase son propios de esa clase 
		\item Los atributos y m\'etodos pueden ser p\'ublicos, privados o heredados
	\end{itemize}
\end{frame}

\section{clases}

\begin{frame}{Clases}
	Son una plantilla sobre la que se construye un objeto, como si del plano de un edificio se tratara o del diseño de un carro, puede resultar algo confuso pero resumiendo, el objeto 
	es cuando ya puedes interactuar con el, en el caso del edificio es cuando ya lo ves construido.
\end{frame}


\subsection{Constructor y destructor}

\begin{frame}{Constructor}
	Es un m\'etodo que  construye el objeto, normalmente se utiliza para darle un estado inicial a los atributos de un objeto pero esto no es obligatorio.
\end{frame}

\begin{frame}{destructor}
	Es un m\'etodo que destruye el objeto, se utiliza para realizar algunas acciones, por ejemplo en el caso de un carro, si el motor falla se enviara la alerta al tablero.
\end{frame}

\subsection{Miembros de una clase}
\begin{frame}{Miembros de una clase}
	Hacen referencia a los componentes de una clase, por ejemplo en un carro, tenemos espejos, motor, llantas, asientos, puertas y m\'as.
	En una clase tendr\'iamos esos mismos componentes pero haciendo referencia mediante un tipo de dato.
\end{frame}

\section{Nivel de acceso}

\begin{frame}{P\'ublico}
	Nos dice que cualquiera puede hacer uso de la clase o de los m\'etodos de la misma, un modo de verlo ser\'ia como tener una casa sin seguro, cualquiera podr\'ia entrar.
\end{frame}

\begin{frame}{Privado}
	Cuando una clase es privada quiere decir que solo la clase puede hacer uso de sus m\'etodos, aqu\'i ser\'ia como tener una casa con seguro y solo quien tenga llave podr\'a entrar
\end{frame}

\begin{frame}{Amigo}
	Aplica cuando queremos que otra clase o una funci\'on tenga acceso a los miembros de una clase, en el caso de la casa, ser\'ia como darle una llave a un amigo o familiar para que pueda entrar.
\end{frame}

\begin{frame}{Protegido}
	Es parecido al nivel amigo, sin embargo aplica para herencia de clases que veremos m\'as adelante, sirve para dar acceso a los miembros privados de una clase a sus clases derivadas, es como cuando tus pap\'as
	te dan permiso para algo.
\end{frame}

\section{Herencia}

\begin{frame}{Clase Padre}
	Es la clase hija principal o de la que las dem\'as heredan y que tiene caracter\'isticas comunes a todas las subclases o clases hijas, un ejemplo ser\'ia una clase veh\'iculo como clase principal y las clases carro y motocicleta
	como subclases 
\end{frame}

\begin{frame}{Clase Hija}
	Es la clase que hereda los m\'etodos y atributos de la clase padre, tienen muchas caracter\'isticas en com\'un pero es necesario diferenciar una subclase de otra, como por ejemplo una motocicleta y un carro, como caracter\'isticas comunes ambos son veh\'iculos a motor y tienen ruedas, pero no son iguales en tamaño ni en capacidad
\end{frame}

\section{Polimorfismo}

\begin{frame}{Polimorfismo}
	Es la capacidad de una clase base de actuar como una de sus clases derivadas, es decir puede adoptar m\'ultiples formas.
\end{frame}

\subsection{Abstracci\'on}

\begin{frame}{Clases abstractas}
	Son clases de las clases no se pueden declarar instancias, es decir solo son una plantilla, ya que solo sirven de base para poder crear las subclases u objeto de sus clases derivadas, un ejemplo ser\'ia
	la clase veh\'iculo, si declaramos una clase veh\'iculo no podr\'iamos saber de que tipo de veh\'iculo se trata.
\end{frame}

\begin{frame}{M\'etodos virtuales}
	Son m\'etodos de una clase abstracta pero que su definici\'on o redefinici\'on se hace en las clases heredadas. Una clase abstracta contiene por lo menos un m\'etodo virtual.
\end{frame}

\begin{frame}{Interfaces}
	Una interfaz es una clase abstracta que solo tiene m\'etodos virtuales, son muy \'utiles para declarar funcionalidades generales .\nocite{BHASIN,BJARNE1,BJARNE2,CAIRO,CPP,DEITEL,DOWNEY,JAWORSKI,KENN,lAAKMANN,MATTHES,RAMALHO,SED,BRASSARD}
\end{frame}

\section{Referencias}
\subsection{}
\begin{frame}[allowframebreaks]
	
	\frametitle{Referencias}
	\bibliographystyle{acm}
	\bibliography{ref.bib}
\end{frame}

\end{document}
